\documentclass[11pt]{article}
\usepackage{fontspec,amsmath,amssymb,amsthm,hyperref}
\usepackage[margin=1in]{geometry}
\setlength{\emergencystretch}{3em}
\title{Exact physical-temperature Gaussian velocity radius}
\author{LangevinLean mathematical exposition}
\date{}
\begin{document}\maketitle
\noindent Source module: \texttt{KIPLMC/Moments.lean}. Compile this proof document with XeLaTeX or LuaLaTeX. The accompanying HTML pairs the mathematical argument with the exact Lean statements and records the scope of the formalization.


In scaled particle coordinates every velocity coordinate has variance \(1/N\). The thermal target is explicitly the image of a standard Gaussian under multiplication by \(N^{-1/2}\). This module proves the reverse normalization as an actual law equality and computes the exact L\ensuremath{{}^{2}} radius \(\sqrt{(d_\theta+Nd_x)/N}\). No independence of coordinates is needed beyond what is already encoded by that Gaussian vector law; zero-dimensional spaces are included.


\section{Undo the thermal normalization as a pushforward equality}


For \(N>0\), the scalar \(\sqrt N\) is nonzero. Composing the target's defining map \(v\mapsto N^{-1/2}v\) with \(v\mapsto\sqrt Nv\) gives the identity on the configuration space. Lean rewrites the two successive pushforwards with \texttt{Measure.map\_map}, proves the two scalar multiplications cancel using \texttt{smul\_smul} and \texttt{mul\_inv\_cancel₀}, then applies \texttt{Measure.map\_id}. Both maps are continuous and hence measurable. This is an exact distributional normalization, rather than a covariance-only assertion.



\noindent Lean declarations: \texttt{velocityTarget\_unscale}.

\section{Read off the L\ensuremath{{}^{2}} moment from the actual Gaussian marginal}


Suppose the L\ensuremath{{}^{2}} random velocity \(V\) has the thermal marginal. The preceding map equality and \texttt{HasLaw.comp} show that \(\sqrt NV\) has the standard Gaussian law. The almost-everywhere scalar-multiplication rule for Lp representatives transfers this statement to the actual Lp class. The shared Gaussian-moment theorem therefore gives \(\|\sqrt NV\|_{L^2}=\sqrt{d_\theta+Nd_x}\). Since \(\sqrt N>0\), the norm scaling rule and division yield

\[\|V\|_{L^2}=\sqrt{\frac{d_\theta+Nd_x}{N}}.\]

Lean simplifies the cardinality of the sum/product coordinate index, rewrites the square root of a quotient, and divides by the strictly positive \(\sqrt N\). The moment is exactly the constant needed by the shared kinetic local-error argument at temperature \(1/N\).



\noindent Lean declarations: \texttt{thermal\_velocity\_Lp\_norm}.

\end{document}
